% Lampiran
\lampiran{1}{Contoh judul lampiran}
Lampiran memuat hasil penelitian yang tidak dimasukkan ke dalam tubuh tulisan agar tidak mengganggu uraian, misalnya tabel yang melebihi satu halaman, peta lokasi, perhitungan, dan prosedur pendukung. Tiap lampiran diberi nomor urut dan judul yang jelas; judul ditulis dengan huruf kecil kecuali huruf awal kalimat dan tanpa titik di akhir.

\lampiran{2}{Hasil pemeriksaan kemiripan}
Hasil pemeriksaan kemiripan (plagiasi) dari perangkat yang ditetapkan Program Studi dilampirkan pada pengajuan ujian. Cantumkan ringkasan hasil seperti Tabel L2.1 beserta tangkapan layar atau laporan dari perangkat tersebut (nama perangkat, persentase, dan tanggal pada tabel adalah contoh).

\begin{singlespace}
\begin{longtable}{|p{0.38\linewidth}|p{0.50\linewidth}|}
\caption{Contoh ringkasan hasil pemeriksaan kemiripan}\\
\hline
Perangkat pemeriksa & Nama perangkat (contoh: Turnitin) \\
\hline
Tanggal pemeriksaan & 15 Mei 2026 \\
\hline
Tingkat kemiripan & 18\% \\
\hline
Batas maksimum Program Studi & 25\% (contoh; sesuai ketentuan Program Studi) \\
\hline
Keterangan & Memenuhi syarat \\
\hline
\end{longtable}
\end{singlespace}
